\section{Finite-window normalization}

We introduce the finite languages and the map $\Fold_m$.  The only
numeration-theoretic input needed later is Zeckendorf uniqueness, which
identifies the golden-mean language as the canonical target of the fold.

\subsection{Admissible words}

Throughout the paper the Fibonacci numbers are normalized by
\[
F_1=F_2=1,
\qquad
F_{k+2}=F_{k+1}+F_k
\quad (k\ge 1).
\]

For $m\ge 1$ and $s\in\{0,1\}$ let
\[
X_m^{(s)}:=\{w=w_1\cdots w_m\in X_m:w_m=s\}.
\]
Thus $X_m^{(0)}$ and $X_m^{(1)}$ split admissible words by their last symbol.

\begin{proposition}[Terminal recursion for admissible words]
\label{prop:stable-syntax-terminal-recursion}
For every $m\ge 1$, if
\[
v_m:=
\begin{pmatrix}
\lvert X_m^{(1)}\rvert\\[2pt]
\lvert X_m^{(0)}\rvert
\end{pmatrix},
\]
then
\[
v_{m+1}=N_\tau v_m,
\qquad
N_\tau=
\begin{pmatrix}
0&1\\[2pt]
1&1
\end{pmatrix}.
\]
\end{proposition}

\begin{proof}
If $w\in X_{m+1}$ ends in $1$, then its penultimate symbol must be $0$.
Deleting the last symbol therefore gives a bijection
\[
X_{m+1}^{(1)}\longleftrightarrow X_m^{(0)},
\]
so $\lvert X_{m+1}^{(1)}\rvert=\lvert X_m^{(0)}\rvert$.

If $w\in X_{m+1}$ ends in $0$, then deleting the last symbol gives an
arbitrary word in $X_m$, and appending $0$ to a word in $X_m$ never creates
the forbidden pattern $11$.  Hence
\[
\lvert X_{m+1}^{(0)}\rvert
=
\lvert X_m\rvert
=
\lvert X_m^{(0)}\rvert+\lvert X_m^{(1)}\rvert.
\]
Combining the two identities gives the displayed matrix recursion.
\end{proof}

\begin{corollary}[Fibonacci counting]
\label{cor:stable-syntax-counting}
For every $m\ge 1$ one has
\[
\lvert X_m\rvert=F_{m+2}.
\]
\end{corollary}

\begin{proof}
The recursion in Proposition~\ref{prop:stable-syntax-terminal-recursion} is
the Fibonacci recursion with the initial vector
\[
v_1=
\begin{pmatrix}
1\\[2pt]
1
\end{pmatrix}.
\]
Therefore
\[
\lvert X_m\rvert
=
\lvert X_m^{(0)}\rvert+\lvert X_m^{(1)}\rvert
=
F_{m+2}.
\qedhere
\]
\end{proof}

\subsection{Fibonacci value and the fold}

Let
\[
\Xfin
:=
\left\{
c=(c_1,c_2,\ldots)\in\{0,1\}^{\NN}:
\supp(c)\text{ finite and }c_kc_{k+1}=0\text{ for all }k
\right\}.
\]
For $c\in\Xfin$ define
\[
\Val(c):=\sum_{k\ge 1} c_kF_{k+1}.
\]
Zeckendorf's theorem gives, for every $N\in\NN$, a unique $Z(N)\in \Xfin$
satisfying $\Val(Z(N))=N$ \cite{Zeckendorf1972}.

For a finite binary window $\omega=\omega_1\cdots\omega_m\in\Omega_m$ define
its Fibonacci value by
\[
N(\omega):=\sum_{k=1}^m \omega_kF_{k+1}.
\]
Let
\[
\pi_m:\Xfin\to X_m,
\qquad
\pi_m(c_1,c_2,\ldots):=(c_1,\ldots,c_m),
\]
be the length-$m$ prefix projection.

\begin{definition}[Finite-window fold]
\label{def:fold-word}
For $m\ge 1$ define
\[
\Fold_m:\Omega_m\to X_m,
\qquad
\Fold_m(\omega):=\pi_m(Z(N(\omega))).
\]
\end{definition}

\begin{definition}[Visible value]
\label{def:visible-value}
For $x=x_1\cdots x_m\in X_m$ define
\[
V_m(x):=\sum_{k=1}^m x_kF_{k+1}.
\]
\end{definition}

\begin{proposition}[Visible-value parametrization of $X_m$]
\label{prop:visible-value-parametrization}
For every $m\ge 1$, the map
\[
V_m:X_m\to \{0,1,\ldots,F_{m+2}-1\}
\]
is a bijection.
\end{proposition}

\begin{proof}
If $x\in X_m$, then
\[
(x_1,\ldots,x_m,0,0,\ldots)\in \Xfin
\]
is already a Zeckendorf digit string, so by uniqueness it is the Zeckendorf
representation of the integer $V_m(x)$.  In particular $V_m(x)<F_{m+2}$,
because otherwise the Zeckendorf representation would have to use the weight
$F_{m+2}$ or a larger one.

Conversely, let $0\le r<F_{m+2}$.  The Zeckendorf representation $Z(r)$ cannot
use the weight $F_{m+2}$ or any larger weight, so it has the form
\[
Z(r)=(x_1,\ldots,x_m,0,0,\ldots)
\]
for a unique $x\in X_m$.  Then $V_m(x)=r$.  Thus $V_m$ is bijective.
\end{proof}

\begin{corollary}[Unique overflow decomposition]
\label{cor:visible-value-overflow}
For every $m\ge 1$ and every $\omega\in\Omega_m$, there exists a unique
$b(\omega)\in\{0,1\}$ such that
\[
N(\omega)=V_m(\Fold_m(\omega))+b(\omega)F_{m+2}.
\]
Equivalently, $\Fold_m(\omega)$ is the unique admissible word in $X_m$ whose
visible value is congruent to $N(\omega)$ modulo $F_{m+2}$.
\end{corollary}

\begin{proof}
Since
\[
0\le N(\omega)\le \sum_{k=1}^m F_{k+1}=F_{m+3}-2<2F_{m+2},
\]
there is a unique $b(\omega)\in\{0,1\}$ such that
\[
r:=N(\omega)-b(\omega)F_{m+2}
\]
lies in $\{0,1,\ldots,F_{m+2}-1\}$.  By
Proposition~\ref{prop:visible-value-parametrization}, there is a unique
$x\in X_m$ with $V_m(x)=r$.  The first $m$ digits of $Z(N(\omega))$ are
precisely this $x$, so $x=\Fold_m(\omega)$.
\end{proof}

For convenience, Table~\ref{tab:value-maps} summarizes the value maps used
throughout the paper.

\begin{table}[t]
\centering
\small
\begin{tabular}{lll}
\toprule
symbol & domain & role \\
\midrule
$N(\omega)$ & $\Omega_m$ & full Fibonacci value of the raw window $\omega$ \\
$V_m(x)$ & $X_m$ & visible value carried by the admissible word $x$ \\
$b(\omega)$ & $\Omega_m$ & overflow bit beyond the visible range $[0,F_{m+2})$ \\
$\Val(a)$ & $\mathcal D$ & value of a finitely supported digit configuration \\
\bottomrule
\end{tabular}
\caption{Value maps used throughout the paper.  The visible value stays within
the length-$m$ window, while the overflow bit records whether one copy of
$F_{m+2}$ has escaped beyond it.}
\label{tab:value-maps}
\end{table}

